\subsection{Metric Applicability for Onocoy}
\label{subsec:onocoy_metric_applicability}

The five standardised metrics defined in Table~\ref{tab:comparative_metrics} are not uniformly applicable to Onocoy. As a capped-supply network with a recent token generation event and near-zero annualised service revenue, several metrics fail to produce interpretable values.

\textbf{30-Day Node Retention} is directly applicable. Provider stickiness is central to GNSS infrastructure sustainability, and retention under stress can be observed from on-chain participation data once a sufficient observation window exists.

\textbf{Burn-to-Mint Ratio} is partially applicable. Onocoy does not use pure Burn-and-Mint Equilibrium; the analogue is Data Credit consumption plus buyback-burn revenue divided by emissions from capped inventory. The resulting ratio is directionally informative but not directly comparable to BME-native networks without noting the mechanism difference.

\textbf{Revenue per Node}, \textbf{Token Turnover}, and \textbf{FDV / Annualised Revenue} are not currently applicable. With negligible service revenue, per-node ratios are undefined. Token Turnover requires sustained secondary-market liquidity, which may be insufficient post-TGE. FDV divided by near-zero revenue produces values that cannot support comparative inference.

These limitations redirect analytical attention toward stress dimensions where Onocoy-specific signals are stronger: competitive yield pressure relative to Geodnet, the subsidy gap between provider costs and realised fiat revenue, elastic provider exit under triple-mining conditions, and sunk-cost retention friction from physical installation requirements.
